\begin{code}[hide]
module vqupit.sections.background where

open import Data.Nat using (ℕ)
open import Data.Product using (_×_ ; _,_ ; Σ-syntax)
open import Relation.Binary.PropositionalEquality using (_≡_ ; _≢_)
open import Notations using (₀)

infixl 7 _•_
infixl 9 _↑
infix 4 _≈ˢ_ _CRel,_===_ _IsPresentationOf_

postulate
  proof-elided : ∀ {ℓ} {A : Set ℓ} → A
  n : ℕ
  Circuit : ℕ → Set
  CZ H S : Circuit n
  _•_ : Circuit n → Circuit n → Circuit n
  _↑ : Circuit n → Circuit n
  _CRel,_===_ : ℕ → Set
  _IsPresentationOf_ : Set → Set → Set
  Pauli⋊Sp : ℕ → Set
  ℤ : ℕ → Set
  ₚ : ℕ
  NF : ℕ → Set
  Word Sp : Set
  [_] : ∀ {n} → NF n → Word
  ⟦_⟧ : Word → Sp
  _≈ˢ_ : Sp → Sp → Set

module PapPres where
  postulate presentation : Set
module U where
  postulate ⟦[]⟧-injective : Set
\end{code}
\section{Qupit Clifford Circuits and the Main Theorems}\label{sec:background}

\subsection{From qubit to qudit Clifford gates}

Fix an odd prime $p$ and let $\omega = e^{2\pi i/p}$.  A \emph{qupit} is
a $p$-dimensional quantum system with computational basis
$\ket{0},\dots,\ket{p-1}$, indices read in $\Zp$.  The single-qupit Pauli
operators are $X\ket{j} = \ket{j+1}$ and $Z\ket{j} = \omega^{j}\ket{j}$;
the $n$-qupit Pauli group $\Pauli$ is generated by tensor products of $X$
and $Z$ together with the phases $\omega^t$.  The $n$-qupit Clifford
group $\Cliff$ is generated, under composition and tensor product, by the
gates of Table~\ref{tab:gates}, whose qubit counterparts are shown for
comparison.  We follow the qupit paper's conventions
\cite{bian2026qudit}: the $\tfrac12$ in $S$ is the inverse of $2$ in
$\Zp$, and $H$ carries a global phase $\delta_p$ chosen so that all three
generators have determinant one (which is what makes the Clifford phases
exactly the powers of $\omega$).

\begin{table}[h]
\caption{The Clifford generators for qubits and for qupits (odd prime $p$).}
\label{tab:gates}
\centering
\begin{tabular}{lll}
\toprule
 & qubit ($p = 2$) & qupit (odd prime $p$) \\
\midrule
$H$  & $\frac{1}{\sqrt2}\begin{pmatrix}1&1\\1&-1\end{pmatrix}$
     & $\ket{j} \mapsto \tfrac{\delta_p}{\sqrt p}\sum_k \omega^{jk}\ket{k}$ \\[2ex]
$S$  & $\begin{pmatrix}1&0\\0&i\end{pmatrix}$
     & $\ket{j} \mapsto \omega^{j(j-1)/2}\,\ket{j}$ \\[2ex]
$CZ$ & $\mathrm{diag}(1,1,1,-1)$
     & $\ket{j}\ket{l} \mapsto \omega^{jl}\,\ket{j}\ket{l}$ \\
\bottomrule
\end{tabular}
\end{table}

Three quotients of $\Cliff$ matter.  Writing $\langle\omega\rangle$ for
the scalars, $\PCliff = \Cliff/\langle\omega\rangle$ is the
\emph{projective} Clifford group and $\PPauli =
\Pauli/\langle\omega\rangle \cong \Zp^{2n}$ the projective Pauli group,
an abelian group in which $X$ and $Z$ commute.  Conjugation by a Clifford
operator permutes Paulis and preserves their commutation phases, so
modulo Paulis a Clifford operator is a symplectic matrix: $\Cliff/\Pauli
\cong \PCliff/\PPauli \cong \Sp$, the symplectic group of the
$2n$-dimensional $\Zp$-vector space $\Zp^{2n}$ with its standard
symplectic form~\cite{bian2026qudit}, and the projective Clifford group
splits: $\PCliff \cong \PPauli \rtimes \Sp$.  These are the ``glue
facts'' of the introduction; our formal development \emph{defines} its
semantic groups as the right-hand sides.

\subsection{Circuits in Agda}

In the formalisation a gate set is a family of gate types indexed by
arity, and a circuit on $n$ wires is a \emph{word} over wire-indexed
generators (\S\ref{sec:framework} gives the generic machinery):

\begin{code}
data SympGate : ℕ → Set where
  H-gate   : SympGate 1
  S-gate   : SympGate 1
  CZ-gate  : SympGate 2

example : Circuit 2
example = CZ • H ↑ • S
\end{code}

$\aid{Circuit}\;n = \aid{Word}\;(\aid{Gen}\;n)$ is the free monoid on
generators; $\aid{{\_}↑}$ shifts a circuit up one wire.  Diagrams are
read in matrix-multiplication order, so \aid{example} is
\fitfig[3cm]{ex-word}.  The remaining usual gates are \emph{words}, not
extra generators: $Z = H H S H H S^{-1}$, $X = H S H H S^{-1} H$, the
swap $\mathit{Ex} = (CZ \cdot H{\downarrow} \cdot H{\uparrow})^{3}$, and
the controlled-$X$ gates $CX = H{\downarrow}^{3} \cdot CZ \cdot
H{\downarrow}$ and $XC$ with the wires exchanged, where $S^{-1} =
S^{p-1}$, $H^{-1} = H^{3}$ and $CZ^{-1} = CZ^{p-1}$ are the words that
the relations make inverses, and \aid{{\_}↓} is the identity on
circuits and exists to make ``the lower copy of a gate'' visible in a
rule.  One more derived gate has no qubit counterpart:
the \emph{multiplier} $M_x\ket{j} = \ket{xj}$ for a unit $x \in
\Zp^{*}$.  The qupit paper spells it (its Lemma~2.14) as $H S^{g^{-1}}
H S^{g} H S^{g^{-1}}$ followed by a Pauli correction and a global
phase; we mostly use $R^{x}\,H\,R^{x^{-1}}H\,R^{x}H$ with $R = S\cdot
Z^{1/2}$, which absorbs the Pauli and drops the phase.  One of the
results below is that the spellings are interderivable.

Everything below lives in one module parameterised by the prime, as
$p = \aid{p-3} + 3$ with a primality proof, and by a primitive root
$g$ of $\Zp^{*}$ with the witness \aid{g-gen} that every unit is a power
of $g$.  So every theorem holds for every odd prime uniformly, and for
the multiplicative generator chosen, exactly as the qupit paper fixes a
generator $g$ in its Figure~1; the existence of a primitive root modulo
every prime is classical and is taken as a parameter rather than proved.

\subsection{Main result 1: a presentation of the projective qupit Clifford group}

Our primary rule set --- the qupit paper's Figure~1 read modulo
scalars, with the multiplier spelled over $R$ and with $(SH)^3 =
\varepsilon$ kept as an axiom --- is, in Agda, one inductive family
\aid{{\_}CRel,{\_}==={\_}} with sixteen constructors, one per rule,
each stated at the smallest number of wires at which it makes sense;
Figure~\ref{fig:rules} shows them as circuit equations, drawn from the
Agda terms, under the constructor names used throughout.  Six are
one-qupit rules: the orders $S^{p} = \varepsilon$, $H^{2} = M_{-1}$ and
$(SH)^{3} = \varepsilon$, the multiplier calculus \aid{M-power}
($M_g^{\,k} = M_{g^k}$) and \aid{semi-MR} ($R\,M_g = M_g\,R^{g^2}$),
and \aid{comm-HHSHHS}.  Seven are two-qupit rules: the orders of $CZ$
and of the swap $\mathit{Ex}$, the commutations of $CZ$ with $S$ and
$M_g$ on the upper wire, the two rules moving $S$ and $H$ across the
swap, and \aid{blake-c12}, which relates two circuits implementing
controlled-$Z$.  Three are three-qupit rules: the Yang--Baxter and
\aid{cz-slide} rules of the swap calculus, and \aid{semi-CX↑-CZ↓}.
Beyond these, \emph{every} circuit theory in
the framework gets the same three structural rules for free: congruence
under $\uparrow$, and gates on disjoint wires commute, one rule per
arity (a global scalar gate, when there is one, gets a fourth: it is
the same gate on every wire; that it then commutes with everything is
a theorem).  Nothing else is an
axiom: $XZ = ZX$ modulo scalars, the commutation of $\mathit{Ex}$ with
$CZ$, and both rules for pushing a Pauli through $CZ$ are \emph{derived}.
The count is sixteen where the paper's Figure~1 has fifteen non-scalar
rules because we keep $(SH)^3 = \varepsilon$ as an axiom under the
$R$-spelling; \S\ref{sec:plumbing} shows that under the paper's own
$S$-spelling it is a theorem, and verifies the paper's fifteen.

\begin{figure}
\centering
\begin{minipage}[t]{0.42\textwidth}\centering\scriptsize
  \begin{tabular}{@{}r@{\;\,}r@{\,}c@{\,}l@{}}
    order-S & \ufig[0.48]{pv0-order-S-l} & $=$ & \ufig[0.48]{pv0-order-S-r}\\[1pt]
    order-H & \ufig[0.48]{pv0-order-H-l} & $=$ & \ufig[0.48]{pv0-order-H-r}\\[1pt]
    M-power & \ufig[0.48]{pv0-M-power-l} & $=$ & \ufig[0.48]{pv0-M-power-r}\\[1pt]
    semi-MR & \ufig[0.48]{pv0-semi-MR-l} & $=$ & \ufig[0.48]{pv0-semi-MR-r}\\[1pt]
    order-CZ & \ufig[0.48]{pv0-order-CZ-l} & $=$ & \ufig[0.48]{pv0-order-CZ-r}\\[1pt]
    order-Ex & \ufig[0.48]{pv0-order-Ex-l} & $=$ & \ufig[0.48]{pv0-order-Ex-r}\\[1pt]
    order-SH & \ufig[0.48]{pv0-order-SH-l} & $=$ & \ufig[0.48]{pv0-order-SH-r}\\
  \end{tabular}
\end{minipage}\hfill
\begin{minipage}[t]{0.55\textwidth}\centering\scriptsize
  \begin{tabular}{@{}r@{\;\,}l@{}}
    comm-CZ-S↑ & \ufig[0.48]{pv0-comm-CZ-Sup}\\[1pt]
    semi-M↑CZ & \ufig[0.48]{pv0-semi-Mup-CZ}\\[1pt]
    semi-Ex-S↑ & \ufig[0.48]{pv0-semi-Ex-Sup}\\[1pt]
    semi-Ex-H↑ & \ufig[0.48]{pv0-semi-Ex-Hup}\\[1pt]
    comm-HHSHHS & \ufig[0.4]{pv0-comm-HHSHHS-l} $=$ \ufig[0.4]{pv0-comm-HHSHHS-r}\\
  \end{tabular}
\end{minipage}\\[3pt]
\scriptsize
\begin{tabular}{@{}r@{\;\,}l@{\qquad}r@{\;\,}l@{}}
  blake-c12 & \ufig[0.48]{pv0-blake-c12} & yang-baxter & \ufig[0.48]{pv0-yang-baxter}\\[2pt]
  cz-slide & \ufig[0.48]{pv0-cz-slide} & semi-CX↑-CZ↓ & \ufig[0.48]{pv0-semi-CXup-CZdn}\\
\end{tabular}
\caption{The sixteen rules of \protect\aid{{\_}CRel,{\_}==={\_}} as
circuit equations.  Diagrams read in matrix-multiplication order, so a
word $w \bullet v$ is drawn with $v$ on the left.
\protect\aid{blake-c12} relates two circuits implementing
controlled-$Z$; \protect\aid{yang-baxter} and \protect\aid{cz-slide}
are the swap calculus.}
\label{fig:rules}
\end{figure}

The first main theorem says that these rules \emph{present} the
projective Clifford group, at every width:

\begin{theorem}[\aid{clifford-presentation}]\label{thm:main}
For every width $n$, the sixteen rules present the projective Clifford
group:
\begin{code}
Theorem-1 : ∀ (n : ℕ) → (n CRel,_===_) IsPresentationOf (Pauli⋊Sp n)
Theorem-1 = PapPres.presentation
\end{code}
\noindent Home: \texttt{ProjectiveClifford.Qupit.Paper-V0.Presentation}
and \texttt{Paper-V1.Presentation}.
\end{theorem}

Here \aid{Pauli⋊Sp}\;$n$ is the semidirect product of the projective
Pauli group $\Zp^{2n}$ by the symplectic group, constructed in
\S\ref{sec:semidirect}, and \aid{{\_}IsPresentationOf{\_}} is the record
of \S\ref{sec:presentations}: a proof that every generator has an
inverse, an interpretation $\sem{\_}$ of words in the group, and a
\emph{group isomorphism} from the words modulo the rules onto the group.
Its two halves are the two directions of the diagram
\[
  \mathsf{Cir}\;n \,/\, {\approx_{\text{16 rules}}}
  \;\xrightarrow{\ \sem{\_}\ \cong\ }\;
  \PPauli \rtimes \Sp \;\cong\; \Cliff/\langle\omega\rangle :
\]
\emph{soundness} --- $\sem{\_}$ is well defined, each rule holds
semantically --- and \emph{completeness} --- $\sem{\_}$ is injective,
semantically equal circuits are related by the rules.  The same theorem
is proved for the paper's fifteen rules in their own spelling
(\S\ref{sec:plumbing}) and, with the scalar $\omega$ and its one
correction added, for the exact Clifford group (\S\ref{sec:scalars}).

\subsection{Main result 2: a unique normal form for symplectic operators}

Symplectic operators are Clifford operators modulo Paulis --- the
``proper'' part of a Clifford operator, its stabiliser tableau.  Our
normal forms for them are \emph{data}, built from four kinds of atomic
boxes:

\begin{minipage}[c]{0.46\textwidth}
\begingroup\renewcommand{\AgdaCodeStyle}{\footnotesize}
\begin{code}
A B D E : Set
A = Σ[ ab ∈ (ℤ ₚ × ℤ ₚ) ] (ab ≢ (₀ , ₀))
B = ℤ ₚ × ℤ ₚ
D = ℤ ₚ × ℤ ₚ
E = ℤ ₚ
\end{code}
\endgroup
\end{minipage}\hfill
\begin{minipage}[c]{0.52\textwidth}\centering\footnotesize
  \begin{tabular}{@{}r@{\;\,}r@{\,}c@{\,}l@{}}
    $A_{a,b}$, $a\neq0$ & \ufig[0.58]{box-A-l} & $=$ & \ufig[0.58]{box-A-r}\\
    $A_{0,b}$, $b\neq0$ & \ufig[0.58]{box-A0-l} & $=$ & \ufig[0.58]{box-A0-r}\\
    $B_{a,b}$, $a\neq0$ & \ufig[0.58]{box-B-l} & $=$ & \ufig[0.58]{box-B-r}\\
    $B_{0,b}$ & \ufig[0.58]{box-B0-l} & $=$ & \ufig[0.58]{box-B0-r}\\
    $D_{a,b}$, $a\neq0$ & \ufig[0.58]{box-D-l} & $=$ & \ufig[0.58]{box-D-r}\\
    $D_{0,b}$ & \ufig[0.58]{box-D0-l} & $=$ & \ufig[0.58]{box-D0-r}\\
    $E_{b}$ & \ufig[0.58]{box-E-l} & $=$ & \ufig[0.58]{box-E-r}\\
  \end{tabular}
\end{minipage}

\medskip\noindent
The boxes are the paper's normal boxes (its Figure~6), with the same
names; the pictures show how each is spelled as a circuit.  Rows of
boxes and the staircase of rows are two more inductive definitions
(\S\ref{sec:nf}), and a normal form is read back as a circuit by
structural recursion.  The second main theorem says the normal form is
unique for the symplectic semantics of \S\ref{sec:symplectic}:

\begin{theorem}[\aid{unique-nf}]\label{thm:unique}
Normal forms with equal symplectic denotations are equal:
\begin{code}
Theorem-2 : ∀ n {u v : NF n} → ⟦ [ u ] ⟧ ≈ˢ ⟦ [ v ] ⟧ → u ≡ v
Theorem-2 = U.⟦[]⟧-injective
\end{code}
\noindent Home: \texttt{Symplectic.Normalization.Uniqueness}.
\end{theorem}

Soundness plus a unique normal form yields completeness, by one generic
lemma (\S\ref{sec:nfproperty}); and the hard, interesting part of the
whole development is to find these coset representatives, parameterised
by the number of wires, and to push every generator through them
(\S\ref{sec:symplectic}).
